% GENERATED FILE. Edit canonical lessons, then run npm run generate:books.
% canonical-source-hash: fnv1a64:7ae4b29392738ca9
% canonical-lessons: ZH-C266-daren, ZH-C266-baobao, ZH-C266-nanhai, ZH-C266-nuhai, ZH-C266-zhongnian

\chapter{Grown-Ups and Children}
\label{ch:zh-grown-ups-and-children}
\hlchaptermodality{266}

\begin{chapteropening}
\textbf{By the end of this chapter:} \emph{I can talk about grown-ups, babies, boys and girls, and middle age.}

Complete the last lesson of chapter 266: I can talk about grown-ups, babies, boys and girls, and middle age.
\end{chapteropening}

\section[\texorpdfstring{\zh{大人} (dàren)}{大人 (dàren)}]{\zh{大人} (dàren) — adult; grown-up}

\label{lesson:ZH-C266-daren}



\begin{quote}
Before the new one: say the Mandarin for stingy, then the Mandarin for nervous.
\end{quote}

\subsection*{You'll want to know: \zh{大人}}
\textbf{\zh{大人}} — \emph{dàren} — \textquotedblleft{}adult; grown-up\textquotedblright{}.

\begin{quote}
\textbf{\zh{大人和小孩都来了}} \emph{dàren hé xiǎohái dōu lái le} — The grown-ups and the children have all come.
\end{quote}

\begin{cousinweb}
\textbf{\zh{大人}} (\emph{dàren}) is built from \zh{大} 'big' and \zh{人} 'a person', both inherited from Middle Chinese, \emph{daH} and \emph{nyin} (Guangyun, in Baxter's transcription). In the word the second loses its tone.
\end{cousinweb}

\begin{grammarlens}[title={What you've built}]
One more word for what I have and can do.
\end{grammarlens}

\subsection*{Your turn}
\begin{itemize}
\raggedright
  \item \emph{Say it:} \emph{dàren}
  \item \emph{Say it:} \emph{dàren}, once more
  \item \emph{Say it:} \emph{dàren hé xiǎohái dōu lái le} — The grown-ups and the children have all come.
  \item \emph{Recall:} say the Mandarin for stingy, then the Mandarin for nervous, then say \emph{dàren} again
\end{itemize}

\begin{tcolorbox}[breakable,colback=teal!4,colframe=teal!35!black,title={Before you move on}]
What does \zh{大人} mean? (\textquotedblleft{}Adult; grown-up\textquotedblright{}.) Say it once more.
\end{tcolorbox}
\section[bǎobao]{bǎobao — baby; darling (said to a small child)}

\label{lesson:ZH-C266-baobao}



\begin{quote}
Before the new one: say the Mandarin for nervous, then the Mandarin for a grown-up.
\end{quote}

\subsection*{You'll want to know: bǎobao}
\emph{bǎobao} — \textquotedblleft{}baby; darling (said to a small child)\textquotedblright{}.

\begin{quote}
\emph{wǒ jiějie de bǎobao hěn kě'ài} — My sister's baby is really cute.
\end{quote}

Its characters come later; for now it is a word for the ear.

\begin{cousinweb}
\emph{bǎobao} is built from \emph{bǎo} 'a treasure', said twice, the second time with no tone. It is inherited from Middle Chinese \emph{pawX} (Guangyun, in Baxter's transcription), glossed 'a treasure'; its use for a baby is not given there.
\end{cousinweb}

\begin{grammarlens}[title={What you've built}]
One more word for what I have and can do.
\end{grammarlens}

\subsection*{Your turn}
\begin{itemize}
\raggedright
  \item \emph{Say it:} \emph{bǎobao}
  \item \emph{Say it:} \emph{bǎobao}, once more
  \item \emph{Say it:} \emph{wǒ jiějie de bǎobao hěn kě'ài} — My sister's baby is really cute.
  \item \emph{Recall:} say the Mandarin for nervous, then the Mandarin for a grown-up, then say \emph{bǎobao} again
\end{itemize}

\begin{tcolorbox}[breakable,colback=teal!4,colframe=teal!35!black,title={Before you move on}]
What does bǎobao mean? (\textquotedblleft{}Baby; darling (said to a small child)\textquotedblright{}.) Say it once more.
\end{tcolorbox}
\section[\texorpdfstring{\zh{男孩} (nánhái)}{男孩 (nánhái)}]{\zh{男孩} (nánhái) — boy}

\label{lesson:ZH-C266-nanhai}



\begin{quote}
Before the new one: say the Mandarin for a grown-up, then the Mandarin for a baby.
\end{quote}

\subsection*{You'll want to know: \zh{男孩}}
\textbf{\zh{男孩}} — \emph{nánhái} — \textquotedblleft{}boy\textquotedblright{}.

\begin{quote}
\textbf{\zh{那个男孩很高}} \emph{nàge nánhái hěn gāo} — That boy is tall.
\end{quote}

\begin{cousinweb}
\textbf{\zh{男孩}} (\emph{nánhái}) is built from \zh{男} 'male' and \zh{孩} 'a child', both inherited from Middle Chinese, \emph{nom} and \emph{hoj} (Guangyun, in Baxter's transcription).
\end{cousinweb}

\begin{grammarlens}[title={What you've built}]
One more word for what I have and can do.
\end{grammarlens}

\subsection*{Your turn}
\begin{itemize}
\raggedright
  \item \emph{Say it:} \emph{nánhái}
  \item \emph{Say it:} \emph{nánhái}, once more
  \item \emph{Say it:} \emph{nàge nánhái hěn gāo} — That boy is tall.
  \item \emph{Recall:} say the Mandarin for a grown-up, then the Mandarin for a baby, then say \emph{nánhái} again
\end{itemize}

\begin{tcolorbox}[breakable,colback=teal!4,colframe=teal!35!black,title={Before you move on}]
What does \zh{男孩} mean? (\textquotedblleft{}Boy\textquotedblright{}.) Say it once more.
\end{tcolorbox}
\section[\texorpdfstring{\zh{女孩} (nǚhái)}{女孩 (nǚhái)}]{\zh{女孩} (nǚhái) — girl}

\label{lesson:ZH-C266-nuhai}



\begin{quote}
Before the new one: say the Mandarin for a baby, then the Mandarin for a boy.
\end{quote}

\subsection*{You'll want to know: \zh{女孩}}
\textbf{\zh{女孩}} — \emph{nǚhái} — \textquotedblleft{}girl\textquotedblright{}.

\begin{quote}
\textbf{\zh{这个女孩很喜欢读书}} \emph{zhège nǚhái hěn xǐhuan dúshū} — This girl loves studying.
\end{quote}

\begin{cousinweb}
\textbf{\zh{女孩}} (\emph{nǚhái}) is built from \zh{女} 'female' and \zh{孩} 'a child', both inherited from Middle Chinese, \emph{nrjoX} and \emph{hoj} (Guangyun, in Baxter's transcription).
\end{cousinweb}

\begin{grammarlens}[title={What you've built}]
One more word for what I have and can do.
\end{grammarlens}

\subsection*{Your turn}
\begin{itemize}
\raggedright
  \item \emph{Say it:} \emph{nǚhái}
  \item \emph{Say it:} \emph{nǚhái}, once more
  \item \emph{Say it:} \emph{zhège nǚhái hěn xǐhuan dúshū} — This girl loves studying.
  \item \emph{Recall:} say the Mandarin for a baby, then the Mandarin for a boy, then say \emph{nǚhái} again
\end{itemize}

\begin{tcolorbox}[breakable,colback=teal!4,colframe=teal!35!black,title={Before you move on}]
What does \zh{女孩} mean? (\textquotedblleft{}Girl\textquotedblright{}.) Say it once more.
\end{tcolorbox}
\section[\texorpdfstring{\zh{中年} (zhōngnián)}{中年 (zhōngnián)}]{\zh{中年} (zhōngnián) — middle age}

\label{lesson:ZH-C266-zhongnian}



\begin{quote}
Before the new one: say the Mandarin for a boy, then the Mandarin for a girl.
\end{quote}

\subsection*{You'll want to know: \zh{中年}}
\textbf{\zh{中年}} — \emph{zhōngnián} — \textquotedblleft{}middle age\textquotedblright{}.

\begin{quote}
\emph{tā bù lǎo shì zhōngnián rén} — He isn't old; he's middle-aged.
\end{quote}

\begin{cousinweb}
\textbf{\zh{中年}} (\emph{zhōngnián}) is built from \zh{中} 'middle' and \zh{年} 'a year', both inherited from Middle Chinese, \emph{trjuwng} and \emph{nen} (Guangyun, in Baxter's transcription).
\end{cousinweb}

\begin{grammarlens}[title={What you've built}]
One more word for what I have and can do.
\end{grammarlens}

\subsection*{Your turn}
\begin{itemize}
\raggedright
  \item \emph{Say it:} \emph{zhōngnián}
  \item \emph{Say it:} \emph{zhōngnián}, once more
  \item \emph{Say it:} \emph{tā bù lǎo shì zhōngnián rén} — He isn't old; he's middle-aged.
  \item \emph{Recall:} say the Mandarin for a boy, then the Mandarin for a girl, then say \emph{zhōngnián} again
\end{itemize}

\begin{tcolorbox}[breakable,colback=teal!4,colframe=teal!35!black,title={Before you move on}]
What does \zh{中年} mean? (\textquotedblleft{}Middle age\textquotedblright{}.) Say it once more.
\end{tcolorbox}
